\documentclass[12pt]{amsart}

%%%%%%%%%This project was adapted from Adam Graham-Squire at High Point University.

\addtolength{\hoffset}{-2.25cm}
\addtolength{\textwidth}{4.5cm}
\addtolength{\voffset}{-2.5cm}
\addtolength{\textheight}{5cm}
\setlength{\parskip}{0pt}
\setlength{\parindent}{0in}

\usepackage{amsthm, amsmath, amssymb}
\usepackage[colorlinks = true, linkcolor = black, citecolor = black, final]{hyperref}
\usepackage{graphicx, multicol}
\usepackage{marvosym, wasysym}

\newcommand{\ds}{\displaystyle}
\DeclareMathOperator{\sech}{sech}



\pagestyle{empty}

\begin{document}

\thispagestyle{empty}

{\scshape MAP 2302} \hfill {\scshape \large Hyperbolic Functions} \hfill {4.3 HW\scshape }
 
\smallskip

\hrule

\bigskip

This project may be completed individually or in a group of two or three students.  If you wish to complete the project as a group, let me know so that I can make the appropriate Blackboard tools available.  You will submit the written report to Project \#1 Assignment by uploading a pdf file.  Please follow all the specifications for a written report project (including those for graphs) that are outlined in the Specifications document.

\bigskip

In this investigative project you'll learn about an important topic, but one that we don't have time to cover together as a class.  You should {\bf not} look up in books or on the internet answers or solutions or explanations of any of the problems or ideas here, except where this is explicitly specified.  You are {\bf investigating} for yourself.  

\medskip

Certain combinations of $e^x$ and $e^{-x}$ are so common in mathematics that we have given them special
names, the {\bf hyperbolic functions}. Although they may look very different, the hyperbolic functions
have many properties, identities, derivatives, and antiderivatives that mimic those of the trigonometric functions,
and that is what we will explore in this project. In particular, we will look at the following
functions, known as the hyperbolic sine, hyperbolic cosine, hyperbolic tangent, and hyperbolic
secant functions:
\begin{multicols}{2}
$$\sinh x = \frac{e^x - e^{-x}}{2}$$

$$\cosh x = \frac{e^x + e^{-x}}{2}$$

$$\tanh x = \frac{\sinh x}{\cosh x}$$

$$\sech x = \frac{1}{\cosh x}$$

\end{multicols}

The reason for the name ``hyperbolic" functions is that they are related to the hyperbola in a similar
way that the trigonometric functions are related to the circle. One application of these functions is
that $\cosh x$ can be used to model a catenary, which is the shape that a heavy flexible cable makes
when it is stretched across a long distance (think of telephone and electrical wires). You will look
at such a graph in question 2.

\bigskip

\bigskip

\begin{enumerate}

\item  ~
\begin{enumerate}
\item  On the set of axes, sketch the graphs of the functions $y = \frac{1}{2}e^x$ and $y = \frac{1}{2}e^{-x}$, for $x$-values between -4 and 4. On that same set of axes, use graphical addition to sketch a graph of the
function $\cosh x = \frac{1}{2}e^x + \frac{1}{2}e^{-x}$. ``Graphical addition" just means that you add the $y$-values together for your two graphs to get a $y$-value for a third graph.
\smallskip
\item  Check the accuracy of your graph by graphing $y = \cosh x$ with a graphing utility.  What are the domain and range of $y = \cosh x$?
\end{enumerate}
\bigskip
\item  Graph four members of the family of curves $y = a\cosh(x/a)$ (Note- this graph is called a
catenary, and should roughly look like a heavy wire stretched between two poles.) That
is, choose four different numbers for $a$ (don't choose zero or one, and at least one of your
choices should be negative) and graph the resulting functions. How does the graph change
as $a$ varies? Explain why this happens.
\bigskip
\item  An even function has reflectional symmetry across the $y$-axis, and an odd function will have
rotational symmetry about the origin. This means that an even function will look like a
mirror image if you flipped it over the $y$-axis and an odd function would look the same if
you spun it 180 degrees around the point $(0,0)$.
\begin{enumerate}
\item  Graph $y = \sinh x$, and $y = \tanh x$. From looking at the graphs, guess whether $\cosh$, $\sinh$ and $\tanh$ are even, odd, or neither.
\smallskip
\item  Prove your assertions in part (a) by using the following definitions: A function $f$ is even if $f(x) = f(-x)$ for all values of $x$.  A function is odd if $f(-x) = -f(x)$ for all values of $x$. For example, the function $f(x) = x^3$ is an odd function because
\begin{align*}
f(-x) & = (-x)^3\\ & = (-x)(-x)(-x) \\ & = -x^3 \\ & = -f(x).\end{align*}
You need to give a similar justification for $\cosh$, $\sinh$ and $\tanh$.
\end{enumerate}
\bigskip
\item  Prove the identity $\cosh^2x - \sinh^2x = 1$.  Start with the left-hand side and use the formulas above and any other valid algebraic manipulations to make the expression look like the right-hand side.
\bigskip
\item  Prove the identity $\sinh(x + y) = \sinh x \cosh y + \cosh x \sinh y$.
\bigskip
\item  Questions 4 and 5 are similar to well-known trigonometric identities. Try to discover another identity for the hyperbolic functions (similar to sum, difference, double or half-angle formulas for trigonometry. If you don't know what these are, try googling ``trigonometric identities". This is an okay thing to look up, since I'm telling you to look it up), and then you must also prove that it is true. {\bf If you are working in a group, you need to discover and prove as many identities as you have group members.}
\bigskip
\item  The differentiation formulas for the hyperbolic functions are similar to those of the trigonometric functions, though the signs are sometimes different.
\begin{enumerate}
\item  Prove that $\ds{\frac{d}{dx}\left(\sinh x\right) = \cosh x}$.  Again, start with the left-hand side, and use formulas and algebra and calculus rules to turn it into the right-hand side.
\smallskip
\item  Find the derivatives of $\cosh x$ and $\tanh x$ given in terms of other hyperbolic trigonometric functions, and prove that they are correct.
\end{enumerate}
\bigskip
\item ~
\begin{enumerate}
\item  We say that a function is one-to-one if it is always increasing or always decreasing. Use calculus to explain why $\sinh$ is a one-to-one function (just saying ``look at the graph" is not enough! You need to talk about slopes, derivatives and their signs).
\item  Find a formula for the derivative of the inverse hyperbolic sine function $y = \sinh^{-1}x$.  You should follow the same method using implicit differentiation that the textbook author uses in section 2.7 (again, you can reference the text here, since I'm telling you to do that) to find the derivative of inverse functions. In order to write the final answer in terms of $x$, use the fact that $\cosh y = \sqrt{1 + \sinh^2y}$.
\smallskip
\item  Use the fact that $\sinh x$ and $\sinh^{-1} x$ are inverse functions to verify that $\sinh^{-1}x = \ln\left(x + \sqrt{x^2 + 1}\right)$.  To do this, you need to show that $\sinh(\sinh^{-1} x) = x$. There may be other ways to verify the equality, but I expect you to prove that $\sinh(\sinh^{-1} x) = x$.
\smallskip
\item  Take the derivative of $\sinh^{-1}x = \ln\left(x + \sqrt{x^2 + 1}\right)$, and show your work. Compare
your answer to what you got in part (b).
\end{enumerate}
\bigskip
\item  Determine the antiderivatives of the four hyperbolic functions, $\sinh x$, $\cosh x$, and $\tanh x$.  You can use formulas and identities you've proven in this project as well as all the techniques of integration we have covered in class.  Clearly show all the steps, simplify the final answers as much as possible, and provide written explanation when necessary.

\end{enumerate}

\end{document}